\section{Gemini Plays Pok\'emon: Additional Evidence}\label{app:gpp}

\cref{fig:gpp_yellow,fig:gpp_battle_complexity} in the main text show Yellow Legacy harness updates and battle-agent structural complexity. This appendix collects the remaining GPP evidence.

\begin{figure}[t]
\centering
\includegraphics[width=\linewidth]{h0_gpp_crystal.pdf}
\caption{Crystal head-to-head comparison under the same harness. (Top) per-500-turn updates on skills and sub-agents for Gemini 2.5 Pro (left) and 3 Pro (right). (Bottom) top-5 most-updated components per model.}
\label{fig:gpp_crystal}
\end{figure}

\begin{table}[t]
\centering
\small
\begin{tabular}{lcc}
\toprule
Opponent & Lifetime attempts & Notes \\
\midrule
Lorelei & 18 & Reached Bruno 15+ times \\
Bruno & 20 & Reached Agatha 12+ times \\
Agatha & 18 & Reached Lance 12+ times \\
Lance & 19 & Reached Champion 4 times \\
Champion Pixel & 4 & Victory on attempt 4 \\
\bottomrule
\end{tabular}
\caption{Yellow Legacy Elite Four lifetime attempt totals. The retries were accompanied by increasingly structured battle prompts (\cref{fig:gpp_battle_complexity}) and persistent written memory, producing a text-encoded decision process across the gauntlet.}
\label{tab:gpp_yellow_e4}
\end{table}

\paragraph{Model comparison on Crystal.} On Pok\'emon Crystal under the same harness, Gemini 3 Pro reached comparable early milestones using roughly half the turns and $\sim 60\%$ fewer tokens than Gemini 2.5 Pro. The largest divergence occurred at Olivine Lighthouse: 3 Pro initially treated the pits as dangerous, then stepped into one after exhausting safer hypotheses and cleared the tower, while 2.5 Pro became trapped in a loop of bad assumptions and spent 16{,}403 turns before obtaining the Fog Badge. \cref{fig:gpp_crystal} shows the corresponding update and fixation patterns.

\paragraph{Failure modes that motivated automation.} GPP also exposed recurring failure modes that the human refiner had to repair between runs: assumptions made without verification (the Goldenrod puzzle ran for days because the agent skipped post-battle NPC dialogue containing the missing hint), brittle tool calls with missing parameters, and limited parallel goal pursuit. These are precisely the failures that \texttt{Continual Harness}'s mid-episode refinement targets.

\subsection{Yellow Legacy Battle-Agent Evolution Checkpoints}\label{app:battle-agent-evolution}

\Cref{fig:gpp_battle_complexity} in the main text plots four graph metrics across the 14 structural checkpoints we extracted from \texttt{custom\_agents.json} over the Elite Four window. For completeness we include the mermaid-rendered decision graph behind each checkpoint across two figures: \cref{fig:yellow-battle-agent-evolution-appendix} shows four canonical checkpoints, and \cref{fig:yellow-battle-agent-evolution-appendix-rest} shows the remaining ten. Every chart uses the same renderer and palette, so structural differences reflect prompt changes, not rendering variance. Node colors encode semantic role: \colorbox[HTML]{CFE4D4}{\texttt{entry}}, \colorbox[HTML]{EEEEEE}{\texttt{analysis}}, \colorbox[HTML]{BED3E5}{\texttt{decision gate}}, \colorbox[HTML]{F4C6AC}{\texttt{terminal action}}.

\begin{figure}[!htbp]
\centering
\begin{minipage}[t]{0.24\linewidth}\centering
\includegraphics[width=\linewidth,height=0.75\textheight,keepaspectratio]{yellow_battle/01_138119_baseline_veto_swarm.pdf}\\[2pt]
{\small \textbf{a1.} Turn 138119: baseline veto swarm}
\end{minipage}\hfill
\begin{minipage}[t]{0.24\linewidth}\centering
\includegraphics[width=\linewidth,height=0.75\textheight,keepaspectratio]{yellow_battle/06_139085_compact_rebuild.pdf}\\[2pt]
{\small \textbf{b1.} Turn 139085: compact rebuild}
\end{minipage}\hfill
\begin{minipage}[t]{0.24\linewidth}\centering
\includegraphics[width=\linewidth,height=0.75\textheight,keepaspectratio]{yellow_battle/10_151441_screen_text_and_prediction.pdf}\\[2pt]
{\small \textbf{c1.} Turn 151441: screen-text + prediction}
\end{minipage}\hfill
\begin{minipage}[t]{0.24\linewidth}\centering
\includegraphics[width=\linewidth,height=0.75\textheight,keepaspectratio]{yellow_battle/12_156631_master_agent_intro.pdf}\\[2pt]
{\small \textbf{d1.} Turn 156631: master-agent intro}
\end{minipage}

\caption{The four Yellow Legacy battle-agent checkpoints marked \textbf{a1/b1/c1/d1} on the complexity plot in \cref{fig:gpp_battle_complexity}. These span the arc from a linear survival-gate chain (a1), through a hard-reset compact rebuild (b1), to a rebuilt-bigger screen-text-grounded program (c1), and finally to a master-agent decomposition (d1) that dispatches to five named sub-checks.}
\label{fig:yellow-battle-agent-evolution-appendix}
\end{figure}

\begin{figure}[!htbp]
\centering
% Row 1
\begin{minipage}[t]{0.24\linewidth}\centering
\includegraphics[width=\linewidth,height=0.45\textheight,keepaspectratio]{yellow_battle/02_138914_level_disparity_veto.pdf}\\[2pt]
{\scriptsize Turn 138914: level-disparity veto}
\end{minipage}\hfill
\begin{minipage}[t]{0.24\linewidth}\centering
\includegraphics[width=\linewidth,height=0.45\textheight,keepaspectratio]{yellow_battle/03_138916_chart_prune_hard_reset.pdf}\\[2pt]
{\scriptsize Turn 138916: chart-prune reset}
\end{minipage}\hfill
\begin{minipage}[t]{0.24\linewidth}\centering
\includegraphics[width=\linewidth,height=0.45\textheight,keepaspectratio]{yellow_battle/04_138923_veto_consolidation.pdf}\\[2pt]
{\scriptsize Turn 138923: veto consolidation}
\end{minipage}\hfill
\begin{minipage}[t]{0.24\linewidth}\centering
\includegraphics[width=\linewidth,height=0.45\textheight,keepaspectratio]{yellow_battle/11_153601_last_stand_strategist.pdf}\\[2pt]
{\scriptsize Turn 153601: last-stand strategist}
\end{minipage}

\vspace{0.8em}

% Row 2
\begin{minipage}[t]{0.24\linewidth}\centering
\includegraphics[width=\linewidth,height=0.25\textheight,keepaspectratio]{yellow_battle/05_138925_minimal_fallback.pdf}\\[2pt]
{\scriptsize Turn 138925: minimal fallback}
\end{minipage}\hfill
\begin{minipage}[t]{0.24\linewidth}\centering
\includegraphics[width=\linewidth,height=0.25\textheight,keepaspectratio]{yellow_battle/07_141323_context_and_viability.pdf}\\[2pt]
{\scriptsize Turn 141323: context + viability}
\end{minipage}\hfill
\begin{minipage}[t]{0.24\linewidth}\centering
\includegraphics[width=\linewidth,height=0.25\textheight,keepaspectratio]{yellow_battle/08_146358_current_opponent_and_hp.pdf}\\[2pt]
{\scriptsize Turn 146358: current opponent + HP}
\end{minipage}\hfill
\begin{minipage}[t]{0.24\linewidth}\centering
\includegraphics[width=\linewidth,height=0.25\textheight,keepaspectratio]{yellow_battle/09_147516_free_turn_and_coverage.pdf}\\[2pt]
{\scriptsize Turn 147516: free turn + coverage}
\end{minipage}

\vspace{0.8em}

% Row 3
\begin{minipage}[t]{0.48\linewidth}\centering
\includegraphics[width=\linewidth,keepaspectratio]{yellow_battle/13_159079_hierarchical_master.pdf}\\[2pt]
{\scriptsize Turn 159079: hierarchical master}
\end{minipage}\hfill
\begin{minipage}[t]{0.48\linewidth}\centering
\includegraphics[width=\linewidth,keepaspectratio]{yellow_battle/14_160511_final_master.pdf}\\[2pt]
{\scriptsize Turn 160511: final master}
\end{minipage}

\caption{The remaining ten Yellow Legacy battle-agent checkpoints, grouped by natural aspect. Row 1: long-chain variants around the first complexity spike and the late ``last-stand'' rewrite. Row 2: medium-hierarchy checkpoints across the rebuild-and-grow-again window. Row 3: the two wide master-agent variants that extend the decomposition introduced at checkpoint 12.}
\label{fig:yellow-battle-agent-evolution-appendix-rest}
\end{figure}

\subsection{Crystal Battle Advisor Evolution Checkpoints}\label{app:crystal-battle-advisor-evolution}

Like the Yellow Legacy Elite Four window, the Crystal run required heavy prompt iteration during its Battle Tower attempt. We extracted 10 structural checkpoints from \texttt{custom\_agents.json} for the \texttt{battle\_advisor} agent. \cref{fig:crystal-battle-advisor-evolution-appendix} shows the first six checkpoints (turns 30k--33k); \cref{fig:crystal-battle-advisor-evolution-appendix-rest} shows the remaining four (turns 33k--36k).

\begin{figure}[!htbp]
\centering
% Row 1
\begin{minipage}[t]{0.32\linewidth}\centering
\includegraphics[width=\linewidth,keepaspectratio]{crystal_battle/01-30242-baseline-matchup-recommender.pdf}\\[2pt]
{\scriptsize Turn 30242: baseline matchup recommender}
\end{minipage}\hfill
\begin{minipage}[t]{0.32\linewidth}\centering
\includegraphics[width=\linewidth,keepaspectratio]{crystal_battle/02-30694-counter-mechanic-hardening.pdf}\\[2pt]
{\scriptsize Turn 30694: counter mechanic hardening}
\end{minipage}\hfill
\begin{minipage}[t]{0.32\linewidth}\centering
\includegraphics[width=\linewidth,keepaspectratio]{crystal_battle/03-31220-battle-tower-legality.pdf}\\[2pt]
{\scriptsize Turn 31220: battle tower legality}
\end{minipage}

\vspace{0.8em}

% Row 2
\begin{minipage}[t]{0.32\linewidth}\centering
\includegraphics[width=\linewidth,keepaspectratio]{crystal_battle/04-31906-fatal-weakness-check.pdf}\\[2pt]
{\scriptsize Turn 31906: fatal weakness check}
\end{minipage}\hfill
\begin{minipage}[t]{0.32\linewidth}\centering
\includegraphics[width=\linewidth,keepaspectratio]{crystal_battle/05-32755-role-based-team-policy.pdf}\\[2pt]
{\scriptsize Turn 32755: role-based team policy}
\end{minipage}\hfill
\begin{minipage}[t]{0.32\linewidth}\centering
\includegraphics[width=\linewidth,keepaspectratio]{crystal_battle/06-33306-survival-check.pdf}\\[2pt]
{\scriptsize Turn 33306: survival check}
\end{minipage}

\caption{Crystal \texttt{battle\_advisor} checkpoints 1--6 from the Battle Tower window. The graphs trace the early evolution from a baseline matchup recommender (1) through legality and weakness-check additions (2--4) to role-based and survival rules (5--6).}
\label{fig:crystal-battle-advisor-evolution-appendix}
\end{figure}

\begin{figure}[!htbp]
\centering
% Row 3 (now its own figure)
\begin{minipage}[t]{0.24\linewidth}\centering
\includegraphics[width=\linewidth,keepaspectratio]{crystal_battle/07-33619-max-stats-and-switch-legality.pdf}\\[2pt]
{\scriptsize Turn 33619: max stats \& switch legality}
\end{minipage}\hfill
\begin{minipage}[t]{0.24\linewidth}\centering
\includegraphics[width=\linewidth,keepaspectratio]{crystal_battle/08-34813-empirical-override-and-paranoia.pdf}\\[2pt]
{\scriptsize Turn 34813: empirical override \& paranoia}
\end{minipage}\hfill
\begin{minipage}[t]{0.24\linewidth}\centering
\includegraphics[width=\linewidth,keepaspectratio]{crystal_battle/09-35466-bahamut-reconfiguration.pdf}\\[2pt]
{\scriptsize Turn 35466: bahamut reconfiguration}
\end{minipage}\hfill
\begin{minipage}[t]{0.24\linewidth}\centering
\includegraphics[width=\linewidth,keepaspectratio]{crystal_battle/10-36217-explicit-speed-calculation.pdf}\\[2pt]
{\scriptsize Turn 36217: explicit speed calculation}
\end{minipage}

\caption{Crystal \texttt{battle\_advisor} checkpoints 7--10. The late-window evolution adds switch-legality bookkeeping (7), an empirical-override / paranoia layer (8), a team-specific reconfiguration (9), and an explicit speed calculation rule (10).}
\label{fig:crystal-battle-advisor-evolution-appendix-rest}
\end{figure}

\subsection{Case Study: The Power Plant Route Loop}\label{app:power-plant-incident}

During the Pok\'emon Yellow Legacy run, the AI agent encountered a 1{,}003-turn stagnation loop on Map ID 4 (Route 4, near Cerulean City). Spanning approximately 3.5 hours on August 29, 2025, this incident provides a documented example of failure modes in runtime tool generation, schema mismatch during prolonged stalls, and the limits of the agent's context horizon for error recovery.

The sequence began when the agent decided its goal was to travel to the Power Plant. To do so, it opened the menu to use the HM ``Fly.'' After repeatedly overshooting its target in the party menu, the agent utilized its tool-generation capabilities to bypass the mechanical inefficiency:
\begin{enumerate}
    \item It deleted its existing \texttt{get\_next\_pokemon\_press} tool.
    \item It wrote a new tool called \texttt{fly\_menu\_navigator}, setting its \texttt{autopress\_buttons} flag to \texttt{true}.
    \item It added a directive to its persistent memory: \emph{``I must use the \texttt{fly\_menu\_navigator} tool as intended and trust its output. The \texttt{get\_next\_pokemon\_press} tool was deleted to make space for \texttt{fly\_menu\_navigator} and should not have been used. This also highlights a failure to immediately use a newly defined tool.''}
\end{enumerate}

\paragraph{Schema Mismatch and Execution Loop.} The agent invoked the tool, passing in the current screen text and its target destination (\texttt{``Power Plant''}). However, the generated tool call did not match the meta-harness execution schema. To execute a tool with \texttt{autopress\_buttons: true}, the agent's JSON output requires the \texttt{buttons\_to\_press} array to be explicitly set to \texttt{["tool"]}. Instead, the agent logically deduced that it needed to scroll down the list of cities, populating the array with \texttt{["Down"]}. Because the meta-harness ignores the \texttt{tools\_to\_call} array if the \texttt{"tool"} button is not explicitly pressed, the system simply pressed \texttt{Down} on the emulator and returned control.

Failing to detect the schema mismatch, the agent recorded in its internal reasoning that the custom tool was successfully executing and returning \texttt{"Down"} as the optimal path. It repeated this exact payload 842 times.

\paragraph{Internal State vs. Execution Reality.} Throughout the 3.5-hour loop, the agent's internal \texttt{thoughts} logs indicated it was actively evaluating its progress, despite the environment remaining largely static.

\begin{quote}
\emph{Turn $\sim$172115:} ``I have completed the mandatory self-assessment... 11. Untested Assumptions \& Confirmation Bias: My assumption about the Power Plant being a Fly destination will be tested by cycling through the menu. I am mindful of confirmation bias. My plan for this turn is to continue navigating the Fly menu... I will use the \texttt{fly\_menu\_navigator} tool to take another step towards 'POWER PLANT'.''
\end{quote}

The agent's internal reasoning indicated it was testing an assumption, but it was caught in a schema error and lacked the necessary game knowledge (the Power Plant is not a valid Fly destination in Generation 1).

\paragraph{Loop Termination.} The loop terminated after 1{,}003 consecutive turns once the agent had manually scrolled through all Kanto cities. When the cursor looped back to ``LAVENDER TOWN'' (its starting point), the agent re-evaluated its state:

\begin{quote}
\emph{Turn $\sim$172322:} ``I am in the Fly menu, and the current highlighted location is LAVENDER TOWN. My goal is to reach the Power Plant. Since Lavender Town is the current map, I have cycled through all available Fly destinations, meaning the Power Plant is not a direct Fly destination. I need to exit the Fly menu and walk to Route 10...''
\end{quote}

The agent stopped calling the custom tool, backed out of the menus, and manually walked north to Route 10.

\paragraph{Observations.}
\begin{itemize}
    \item \textbf{Context Horizon Limits:} Tool generation primarily occurred within the first 50-200 turns of encountering a new bottleneck. Beyond 500 turns in a stall, the agent ceased tool creation and reverted to repeated execution patterns.
    \item \textbf{Schema Fragility:} When enforcing new rules via persistent memory, the model remained susceptible to schema mismatch, executing actions that did not align with its intended tool design.
    \item \textbf{Feedback Blindness:} The assumption that the new tool was functioning correctly caused the agent to ignore standard environmental feedback and anomaly detection mechanisms for an extended period.
\end{itemize}
